\section{Arquitectura de soluciones: patrón de colaboración entre dos agentes}

La solución propuesta por Anthropic incluye dos roles clave: Initializer Agent y Coding Agent. Ambos intercambian información a través del estado compartido del sistema de archivos (específicamente los archivos \texttt{progress} y \texttt{git history}), permitiendo la transferencia de datos entre sesiones.

\subsection{Initializer Agent: configuración del entorno}

El Initializer Agent solo se ejecuta durante la primera sesión del proyecto, encargándose de crear el entorno necesario para todas las sesiones posteriores de Coding Agent. Sus responsabilidades específicas incluyen:

\begin{enumerate}
    \item Crear el script \texttt{init.sh}, utilizado para iniciar servidores de desarrollo e infraestructura similar
    \item Crear el archivo \texttt{claude-progress.txt}, que registra los registros de trabajo de cada sesión de agente
    \item Generar un commit inicial de git, registrando todos los archivos añadidos
    \item Redactar la lista de funcionalidades (feature list), expandiendo una lista completa de requisitos basada en el prompt del usuario
\end{enumerate}

\subsection{Coding Agent: avance incremental}

Cada sesión subsiguiente ejecuta al Coding Agent, cuyo comportamiento sigue este patrón:

\begin{enumerate}
    \item Leer los archivos \texttt{progress} y el registro de git (\texttt{git log}) para comprender el progreso actual
    \item Seleccionar la siguiente funcionalidad pendiente de implementación a partir de la feature list
    \item Implementar dicha funcionalidad y realizar pruebas
    \item Realizar un commit en git y actualizar el archivo \texttt{progress}
    \item Terminar la sesión una vez asegurado que el entorno se encuentra en estado limpio
\end{enumerate}

\begin{knowledgebox}{¿Por qué utilizar dos agentes distintos?}
Diferenciar la "ejecución inicial" de la "ejecución posterior" mediante prompts distintos es una elección de diseño importante. El prompt del Initializer Agent se centra en la configuración del entorno y el análisis de requisitos, mientras que el prompt del Coding Agent se enfoca en el desarrollo incremental y la calidad del código. Esta división evita que un único prompt sea excesivamente complejo y permite que cada agente se concentre mejor en sus responsabilidades respectivas. La guía de prompting para Claude 4 de Anthropic también recomienda esta estructura de harness donde "la primera ventana de contexto utiliza un prompt diferente".
\end{knowledgebox}

\subsection{Mecanismo de transferencia de información}

La transferencia de información entre ambos agentes depende completamente del sistema de archivos, y no de ninguna memoria o base de datos:

\begin{table}[H]
\centering
\begin{tabular}{p{4cm}p{4cm}p{5cm}}
\toprule
\textbf{Portador de información} & \textbf{Función} & \textbf{Puntos clave del diseño} \\
\midrule
\texttt{claude-progress.txt} & Registra un resumen del trabajo realizado en cada sesión & Formato de texto libre, facilitando que el agente comprenda rápidamente el contexto \\
\midrule
\texttt{feature\_list.json} & Rastrea el estado de finalización de todas las funcionalidades & Formato JSON, impidiendo que el agente modifique arbitrariamente el contenido \\
\midrule
Historial de git (Git history) & Registra todos los cambios en el código & Mensajes descriptivos para el commit, admitiendo operaciones de revert \\
\midrule
\texttt{init.sh} & Script de inicio del entorno & El agente no necesita rediscover cómo ejecutar el proyecto \\
\bottomrule
\end{tabular}
\caption{Cuatro portadores para la transferencia de información entre sesiones}
\end{table}

\subsection{Resumen de la sección}

La arquitectura de dos agentes descompone el objetivo poco realista de "completar todo en una sola vez" en "configurar el entorno + avanzar de manera incremental continua". El Initializer Agent se encarga de crear toda la infraestructura necesaria, mientras que el Coding Agent realiza avances incrementales manejables en cada sesión. Ambos intercambian estado a través del archivo \texttt{progress}, la feature list, el historial de git y los scripts de inicio.

